% F085｜八类模型速览：保留原 statementbox 框架。
\FReserve{4cm}
\begin{statementbox}
	\textbf{\textcolor{CStatement}{为什么要学这八类模型？}}

	抽象函数题只给 $f(x+y)$、$f(xy)$ 等关系，却不告诉你 $f(x)$ 长什么样。\textbf{不必先猜出解析式！}先观察“输入端”怎么运算，再看“输出端”怎么运算，就可能直接求值。

	例如 $f(x+y)=f(x)f(y)$，若 $f(1)=3$，那么
	\[
		f(3)=f(1+1+1)=[f(1)]^3=27.
	\]
	这正是\textbf{指数型}关系：自变量相加，函数值相乘。

	% F085｜真实数学图；代表函数不构成抽象解的唯一性证明。
\begin{center}
\begin{minipage}{.49\linewidth}\centering
\begin{tikzpicture}
\begin{axis}[width=6.1cm,height=4.8cm,axis lines=middle,
 xmin=-1.3,xmax=3.5,ymin=-0.5,ymax=9,
 xtick={0,1,2,3},xticklabels=\empty,ytick={0,2,4,8},xlabel={$t$},ylabel={$f(t)$},
 title={加变乘：$2^t$},title style={font=\small,yshift=3pt},
 tick label style={font=\small},label style={font=\small},
 ylabel style={at={(axis description cs:0,1)},anchor=south east,xshift=-18pt},
 clip=false,grid=major,grid style={black!8},enlargelimits=false,
 every axis plot/.append style={thick},samples=161]
\addplot[CExplanation,very thick,domain=-1:3] {2^x};
\draw[dashed,black!35] (axis cs:1,0) |- (axis cs:0,2);
\draw[dashed,black!35] (axis cs:2,0) |- (axis cs:0,4);
\draw[dashed,black!35] (axis cs:3,0) |- (axis cs:0,8);
\addplot[only marks,mark=*,mark size=2pt,orange!80!black] coordinates {(1,2) (2,4) (3,8)};

\node[font=\small,anchor=north,yshift=-3pt] at (axis cs:1,-0.5) {$1$};
\node[font=\small,anchor=north,yshift=-3pt] at (axis cs:2,-0.5) {$2$};
\node[font=\small,anchor=north,yshift=-3pt] at (axis cs:3,-0.5) {$3$};

\end{axis}
\end{tikzpicture}
\end{minipage}\hfill
\begin{minipage}{.49\linewidth}\centering
\begin{tikzpicture}
\begin{axis}[width=6.1cm,height=4.8cm,axis lines=middle,
 xmin=-0.3,xmax=8.8,ymin=-2,ymax=3.7,
 xtick={1,2,4,8},xticklabels=\empty,ytick={-1,0,1,2,3},xlabel={$t$},ylabel={$f(t)$},
 title={乘变加：$\log_2t$},title style={font=\small,yshift=3pt},
 tick label style={font=\small},label style={font=\small},
 ylabel style={at={(axis description cs:0,1)},anchor=south east,xshift=-18pt},
 clip=false,grid=major,grid style={black!8},enlargelimits=false,
 every axis plot/.append style={thick},samples=161]
\addplot[CExplanation,very thick,domain=.3:8] {ln(x)/ln(2)};
\draw[dashed,black!35] (axis cs:2,0) |- (axis cs:0,1);
\draw[dashed,black!35] (axis cs:4,0) |- (axis cs:0,2);
\draw[dashed,black!35] (axis cs:8,0) |- (axis cs:0,3);
\addplot[only marks,mark=*,mark size=2pt,orange!80!black] coordinates {(2,1) (4,2) (8,3)};

\node[font=\small,anchor=north,yshift=-3pt] at (axis cs:1,-2) {$1$};
\node[font=\small,anchor=north,yshift=-3pt] at (axis cs:2,-2) {$2$};
\node[font=\small,anchor=north,yshift=-3pt] at (axis cs:4,-2) {$4$};
\node[font=\small,anchor=north,yshift=-3pt] at (axis cs:8,-2) {$8$};

\end{axis}
\end{tikzpicture}
\end{minipage}
\par\small 同一组数反向阅读：$2^{1+2}=2\cdot4=8$；$\log_2(2\cdot4)=1+2=3$。
\end{center}


	\textbf{一眼识别：}加变乘想指数；乘变加想对数；乘变乘想幂；加变加想正比例；和差相加想余弦；和差相乘想正弦；分式和差想正切、余切。
\end{statementbox}

\FReserve{4cm}
\begin{statementbox}
	\textbf{\textcolor{CStatement}{四个代数模型（第 1—4 类）}}
	\[
		\begin{array}{cll}
			\text{1. 指数型}  & f(x+y)=f(x)f(y)  & a^x      \\[4pt]
			\text{2. 对数型}  & f(xy)=f(x)+f(y)  & \log_a x \\[4pt]
			\text{3. 幂函数型} & f(xy)=f(x)f(y)   & x^a      \\[4pt]
			\text{4. 正比例型} & f(x+y)=f(x)+f(y) & kx
		\end{array}
	\]

	\textbf{先记清定义域：}1、4 通常以 $x,y\in\mathbb R$ 为例；2、3 通常先在 $x,y>0$ 讨论；$a^x$ 中 $a>0$，对数还要求 $a\ne1$。每个模型都要求涉及的输入有定义。

	\textbf{最实用的对比：}
	\[
		\underbrace{f(x^n)=n f(x)}_{\text{对数型}},\qquad
		\underbrace{f(x^n)=[f(x)]^n}_{\text{幂函数型}}\quad(n\in\mathbb Z).
	\]
	这里 $x>0$，幂型取不恒为零的解；前者“次数提到前面”，后者“函数值整体乘方”。
\end{statementbox}

\FReserve{4cm}
\begin{statementbox}
	\textbf{\textcolor{CStatement}{四个三角函数模型（第 5—8 类）}}
	\[
		\begin{aligned}
			\text{5. 余弦型}\quad & f(x+y)+f(x-y)=2f(x)f(y),             \\[4pt]
			\text{6. 正切型}\quad & f(x+y)=\frac{f(x)+f(y)}{1-f(x)f(y)}, \\[5pt]
			\text{7. 正弦型}\quad & f(x+y)f(x-y)=f^2(x)-f^2(y),          \\[4pt]
			\text{8. 余切型}\quad & f(x+y)=\frac{f(x)f(y)-1}{f(x)+f(y)}.
		\end{aligned}
	\]
	典型例子依次为 $\cos x$、$\tan x$、$\sin x$、$\cot x$。

	\textbf{分式模型一定先检查：}正切型的 $1-f(x)f(y)\ne0$；余切型的 $f(x)+f(y)\ne0$，而且 $x,y,x+y$ 必须都属于相应定义域。

	\textbf{最重要的提醒：}以上是\textbf{识别线索}，不是无条件的“唯一解析式”。如 $f(x)=x$ 也满足7；5还允许 $f(x)\equiv1$。只有增加足够的条件，才可能确定具体解析式。
\end{statementbox}

